\section{AIME II 2003}
        \begin{enumerate}
        	\item (\textit{AIME II 2003}) The product $N$ of three positive integers is 6 times their sum, and one of the integers is the sum of the other two. Find the sum of all possible values of $N$.


 
 \textbf{Answer: }336



\textbf{Solution}

Let the three integers be $a, b, c$.  $N = abc = 6(a + b + c)$ and $c = a + b$.  Then $N = ab(a + b) = 6(a + b + a + b) = 12(a + b)$.  Since $a$ and $b$ are positive, $ab = 12$ so $\{a, b\}$ is one of $\{1, 12\}, \{2, 6\}, \{3, 4\}$ so $a + b$ is one of $13, 8, 7$ so the sum of all possible values of $N$ is $12 \cdot (13 + 8 + 7) = 12(28) = \boxed{336}$.

 




	\item (\textit{AIME II 2003}) Let $N$ be the greatest integer multiple of 8, whose digits are all different. What is the remainder when $N$ is divided by 1000?


 
 \textbf{Answer: }120



\textbf{Solution}

We want a number with no digits repeating, so we can only use the digits $0-9$ once in constructing our number. To make the greatest number, we want the greatest digit to occupy the leftmost side and the least digit to occupy the rightmost side. Therefore, the last three digits of the greatest number should be an arrangement of the digits $0,1,2$. Since the number has to be divisible by 8, the integer formed by the arrangement of $0,1,2$ is also divisible by 8. The only arrangement that works is $120$. 

 Therefore, the remainder when the number is divided by $1000$ is $\boxed{120}$.

 Note: The number is $9876543120$

 





\textbf{Solution 2}

We notice that the number doesn't matter after the thousands place because any place higher will be divisible by 8. Therefore, the number will look like this(with the maximizing constraint):

 $9876543abc$

 The abc is there because we don't have the guarantee of divisibility by 8. abc will be a rearrangement of $0, 1, 2$

 Now all we have to do is look for the smallest rearrangement of 0, 1, 2. We can write a congruence statement as shown: 

 $100a+10b+10c \cong 0 \mod 8$

 which simplifies to:

 $4a+2b+c \cong 0 \mod 8$

 Now, since we want to minimize the number, we can first try the case when a = 0. 

 This doesn't work, as 2+0 isn't congruent to 0 mod 8. 

 Now lets try a = 1.

 We have $4 + 2b + c \cong 0 \mod 8$

 From which we get 120 as the smallest. 

 Now, taking adding that to the number, we get 9876543120. Getting the remainder mod 1000 gives $\boxed{120}$

 




	\item (\textit{AIME II 2003}) Define a $\text{good word}$ as a sequence of letters that consists only of the letters $A$, $B$, and $C$ - some of these letters may not appear in the sequence - and in which $A$ is never immediately followed by $B$, $B$ is never immediately followed by $C$, and $C$ is never immediately followed by $A$. How many seven-letter good words are there?


 
 \textbf{Answer: }192



\textbf{Solution 1}

There are three letters to make the first letter in the sequence. However, after the first letter (whatever it is), only two letters can follow it, since one of the letters is restricted. Therefore, the number of seven-letter good words is $3*2^6=192$

 Therefore, there are $\boxed{192}$ seven-letter good words.

 





\textbf{Solution 2 (Recursion)}

We solve this problem using recursion. Let $f(x)$ be the number of $x$-letter good words. Thus $f(1) = 3$ (A, B or C) and the answer is just $f(7)$.  The recurrence relation can be found by considering the last letter of one of the valid strings of length $x - 1$.  There are $2$ possibilities for the next letter and thus $f(x) = 2 \cdot f(x-1)$.  Now we can find a closed form as $f(x) = 3 \cdot 2 ^{x-1}$  (easy to prove by induction) and thus $f(7) = 64 * 3 = \boxed{192}$ seven-letter good words.  ~AK2006

 





\textbf{Solution 3}

This restriction forces the letters to appear in $\emph{blocks}$ in a fixed cyclic order: 
 \[A \;\to\; C \;\to\; B \;\to\; A \;\to\; C \;\to\; B \;\to \dots\]For example, if a word starts with  $A$, $\text{ it must look like: }$$\underbrace{AAA\cdots}_{\text{A-block}} \;  \underbrace{CCC\cdots}_{\text{C-block}} \;  \underbrace{BBB\cdots}_{\text{B-block}} \;  \underbrace{AAA\cdots}_{\text{A-block}} \; \dots$

 The word's structure is therefore determined by the starting letter and where the block transitions occur. In a $7$ letter word, there are $6$ gaps between the letters. At each gap, we can either stay in the block or move on to the next. This gives $2$ choices for each gap, or $2^6$ possibilities. Multiply this by the $3$ possible starting letters to get $3*2^6 = \boxed{192}$.

 




	\item (\textit{AIME II 2003}) In a regular tetrahedron, the centers of the four faces are the vertices of a smaller tetrahedron. The ratio of the volume of the smaller tetrahedron to that of the larger is $m/n$, where $m$ and $n$ are relatively prime positive integers. Find $m+n$.


 
 \textbf{Answer: }028



\textbf{Solution}

\textbf{Official Solution (MAA)} Let $O$ and $P$ be the centers of faces $DAB$ and $ABC$, respectively, of regular tetrahedron $ABCD$. Both $\overrightarrow{DO}$ and $\overrightarrow{CP}$ intersect $AB$ at its midpoint $M$. Since $\frac{MO}{MD}=\frac{MP}{MC}=\frac{1}{3}$, triangles $MOP$ and $MDC$ are similar, and $OP = (1/3)DC$. Because the tetrahedra are similar, the ratio of their volumes is the cube of the ratio of a pair of corresponding sides, namely, $(1/3)^3 = 1/27$, so $m + n = \boxed{28}$.
                    \begin{figure}[H]
                        \centering
                        \includegraphics[width=0.5\textwidth]{Solution/AIME_2003_II_4_Solution_Figure_1.png}
                    \end{figure}
                    

 





\textbf{Solution 1}

Embed the tetrahedron in 4-space to make calculations easier.Its vertices are$(1,0,0,0)$, $(0,1,0,0)$, $(0,0,1,0)$, $(0,0,0,1)$.

 To get the center of any face, we take the average of the three coordinates of that face. The vertices of the center of the faces are:$(\frac{1}{3}, \frac{1}{3}, \frac{1}{3}, 0)$,$(\frac{1}{3}, \frac{1}{3},0, \frac{1}{3})$,$(\frac{1}{3},0, \frac{1}{3}, \frac{1}{3})$,$(0,\frac{1}{3}, \frac{1}{3}, \frac{1}{3})$.

 The side length of the large tetrahedron is $\sqrt{2}$ by the distance formula.The side length of the smaller tetrahedron is $\frac{\sqrt{2}}{3}$ by the distance formula.

 Their ratio is $1:3$, so the ratio of their volumes is $\left(\frac{1}{3}\right)^3 = \frac{1}{27}$.

 $m+n = 1 + 27 = \boxed{28}$.

 





\textbf{Solution 2}

Let the large tetrahedron be $ABCD$, and the small tetrahedron be $WXYZ$, with $W$ on $ABC$, $X$ on $BCD$, $Y$ on $ACD$, and $Z$ on $ABD$. Clearly, the two regular tetrahedrons are similar, so if we can find the ratio of the sides, we can find the ratio of the volumes. Let $AB=1$, for our convenience. Dropping an altitude from $W$ to $BC$, and calling the foot $M$, we have $WM=XM=\frac{\sqrt3}{6}$. Since $\cos\angle{WMX}=\cos\angle{AMX}=MX/AM=1/3$. By Law of Cosines, we have $WX=\sqrt{1/12+1/12-2(1/12)(1/3)}=1/3$. Hence, the ratio of the volumes is $\left(\frac{1}{3}\right)^3=1/27$. $m+n=1+27=\boxed{028}$

 





\textbf{Solution 3}

Consider the large tetrahedron $ABCD$ and the smaller tetrahedron $WXYZ$.  Label the points as you wish, but dropping an altitude from the top vertex of $ABCD$, we see it hits the center of the base face of $ABCD$.  This center is also one vertex of $WXYZ$.  Consider a "side" face of $ABCD$, and the center of that face, which is another vertex of $WXYZ$.  Draw the altitude of this side face (which is an equilateral triangle).  These two altitudes form a right triangle.  Since the center of the Side face splits the altitude of the side face into segments in the ratio of $2:1$ (centroid), and since the bases of $WXYZ$ and $ABCD$ are parallel, we can say that the altitudes of tetrahedron $ABCD$ and $WXYZ$ are in the ratio $3:1$.  Thus we compute $\left(\frac{1}{3}\right)^3$, and find $\frac{1}{27}$.  The sum of the numerator and denominator is thus $28$.

 \textbf{Video Solution by Sal Khan} - AMBRIGGS

 




	\item (\textit{AIME II 2003}) A cylindrical log has diameter $12$ inches. A wedge is cut from the log by making two planar cuts that go entirely through the log. The first is perpendicular to the axis of the cylinder, and the plane of the second cut forms a $45^\circ$ angle with the plane of the first cut. The intersection of these two planes has exactly one point in common with the log. The number of cubic inches in the wedge can be expressed as $n\pi$, where n is a positive integer. Find $n$.


 
 \textbf{Answer: }216



\textbf{Solution}

The volume of the wedge is half the volume of a cylinder with height $12$ and radius $6$. (Imagine taking another identical wedge and sticking it to the existing one).  Thus, $V=\dfrac{6^2\cdot 12\pi}{2}=216\pi$, so $n=\boxed{216}$.

 - AMBRIGGS

 




	\item (\textit{AIME II 2003}) In triangle $ABC,$ $AB = 13,$ $BC = 14,$ $AC = 15,$ and point $G$ is the intersection of the medians. Points $A',$ $B',$ and $C',$ are the images of $A,$ $B,$ and $C,$ respectively, after a $180^\circ$ rotation about $G.$ What is the area of the union of the two regions enclosed by the triangles $ABC$ and $A'B'C'?$


 
 \textbf{Answer: }112



\textbf{Solution}

Since a $13-14-15$ triangle is a $5-12-13$ triangle and a $9-12-15$ triangle "glued" together on the $12$ side, $[ABC]=\frac{1}{2}\cdot12\cdot14=84$. 

 There are six points of intersection between $\Delta ABC$ and $\Delta A'B'C'$. Connect each of these points to $G$. 

 
                    \begin{figure}[H]
                        \centering
                        \includegraphics[width=0.5\textwidth]{Solution/AIME_2003_II_6_Solution_Figure_1.png}
                    \end{figure}
                    

 There are $12$ smaller congruent triangles which make up the desired area. Also, $\Delta ABC$ is made up of $9$ of such triangles. Therefore, $\left[\Delta ABC \bigcup \Delta A'B'C'\right] = \frac{12}{9}[\Delta ABC]= \frac{4}{3}\cdot84=\boxed{112}$.

 





\textbf{Solution 2(Doesn’t require good diagram)}

First, find the area of $\Delta ABC$ either like the first solution or by using Heron’s Formula. Then, draw the medians from $G$ to each of $A, B, C, A’, B’,$ and $C’$. Since the medians of a triangle divide the triangle into 6 triangles with equal area, we can find that each of the 6 outer triangles have equal area. (Proof: Since I’m too lazy to draw out a diagram, I’ll just have you borrow the one above. Draw medians $GA$ and $GB’$, and let’s call the points that $GA$ intersects $C’B’$ “$H$” and the point $GB’$ intersects $AC$ “$I$”. From the previous property and the fact that both $\Delta ABC$ and $\Delta A’B’C’$ are congruent, $\Delta GHB’$ has the same area as $\Delta GIA$. Because of that, both “half” triangles created also have the same area. The same logic can be applied to all other triangles).Also, since the centroid of a triangle divides each median with the ratio $2:1$, along with the previous fact, each outer triangle has $1/9$ the area of $\Delta ABC$ and $\Delta A’B’C’$. Thus, the area of the region required is $\frac{4}{3}$ times the area of $\Delta ABC$ which is $\boxed {112}$.

 Solution by Someonenumber011

 





\textbf{Solution 3 (Rigorous)}

$[ABC]$ can be calculated as 84 using Heron's formula or other methods. Since a $180^{\circ}$ rotation is equivalent to reflection through a point, we have a homothety with scale factor $-1$ from $ABC$ to $A'B'C'$ through the centroid $G$. Let $M$ be the midpoint of $BC$ which maps to $M'$ and note that $A'G=AG=2GM,$ implying that $GM=MA'.$ Similarly, we have $AM'=M'G=GM=MA'.$ Also let $D$ and $E$ be the intersections of $BC$ with $A'B'$ and $A'C',$ respectively. The homothety implies that we must have $DE || B'C',$ so there is in fact another homothety centered at $A'$ taking $A'DE$ to $A'B'C'$. Since $A'M'=3A'M,$ the scale factor of this homothety is 3 and thus $[A'DE]=\frac{1}{9}[A'B'C']=\frac{1}{9}[ABC].$ We can apply similar reasoning to the other small triangles in $A'B'C'$ not contained within $ABC$, so our final answer is $[ABC]+3\cdot\frac{1}{9}[ABC]=\boxed{112.}$

 





\textbf{Solution 4}

We determine the area of $\triangle ABC$ to be 84 via Heron's Formulae. A "$180^\circ$ rotation" is the same as just reflecting the figure by centroid $G$. Also, a rotation (or reflection by $G$), simply forms a new triangle that is exactly congruent through SSS to the original. We have established $ABC$ and $A'B'C'$ are congruent.

 Now, we show that the partitioning from the consecutive cuts of the two triangles results in twelve congruent triangles (three triangles formed by each line. Each line, $AB$, $A'B'$, so on so forth are parallel, forming transversals. Repeated parallel base theorems give the desired result.

 Exactly nine of these congruent triangles form the interior of the triangle $ABC$. Say these all have area $T$. Then, $9T = 84$ so $T = 28/3$. Then, multiplying by the 12 triangles, the answer is $28 \cdot 4 = \boxed{112}$.

 




	\item (\textit{AIME II 2003}) Find the area of rhombus $ABCD$ given that the radii of the circles circumscribed around triangles $ABD$ and $ACD$ are $12.5$ and $25$, respectively.


 
 \textbf{Answer: }400



\textbf{Solution}

The diagonals of the rhombus perpendicularly bisect each other. Call half of diagonal BD $a$ and half of diagonal AC $b$. The length of the four sides of the rhombus is $\sqrt{a^2+b^2}$. 

 The area of any triangle can be expressed as $\frac{a\cdot b\cdot c}{4R}$, where $a$, $b$, and $c$ are the sides and $R$ is the circumradius. Thus, the area of $\triangle ABD$ is $ab=2a(a^2+b^2)/(4\cdot12.5)$. Also, the area of $\triangle ABC$ is $ab=2b(a^2+b^2)/(4\cdot25)$. Setting these two expressions equal to each other and simplifying gives $b=2a$. Substitution yields $a=10$ and $b=20$, so the area of the rhombus is $20\cdot40/2=\boxed{400}$.

 





\textbf{Solution 2}

Let $\theta=\angle BDA$. Let $AB=BC=CD=x$. By the extended law of sines,
 \[\frac{x}{\sin\theta}=25\]Since $AC\perp BD$, $\angle CAD=90-\theta$, so
 \[\frac{x}{\sin(90-\theta)=\cos\theta}=50\]Hence $x=25\sin\theta=50\cos\theta$. Solving $\tan\theta=2$, $\sin\theta=\frac{2}{\sqrt{5}}, \cos\theta=\frac{1}{\sqrt{5}}$. Thus
 \[x=25\frac{2}{\sqrt{5}}\implies x^2=500\]The height of the rhombus is $x\sin(2\theta)=2x\sin\theta\cos\theta$, so we want
 \[2x^2\sin\theta\cos\theta=\boxed{400}\]

 




	\item (\textit{AIME II 2003}) Find the eighth term of the sequence $1440,$ $1716,$ $1848,\ldots,$ whose terms are formed by multiplying the corresponding terms of two arithmetic sequences.


 
 \textbf{Answer: }348



\textbf{Solution 0.67}

If you multiply the corresponding terms of two arithmetic sequences, you get the terms of a quadratic function. Thus, we have a quadratic $ax^2+bx+c$ such that $f(1)=1440$, $f(2)=1716$, and $f(3)=1848$. Plugging in the values for x gives us a system of three equations:

 $a+b+c=1440$

 $4a+2b+c=1716$

 $9a+3b+c=1848$

 Solving gives $a=-72, b=492,$ and $c=1020$. Thus, the answer is $-72(8)^2+492\cdot8+1020= \boxed{348}.$

 
Note: WHO WROTE THIS?! ITS THE BIG '26!!~Aarav22

 \textbf{Rigorous Justification for Quadratic} Define $c_n = a_n \cdot b_n$.

 Through definition, $a_n = a + (n-1)d_1$ where $a$ is the first term of $a_n$. Define $b_n$ analogously with difference $d_2$. Therefore,

 
 \[c_n = ab + a(n - 1)d_2 + b(n - 1)d_1 + (n - 1)^2 d_1d_2\]

 Expanding and rearranging in terms of $n$, we have

 
 \[c_n = (d_1d_2)n^2 + (ad_2 + bd_1 - 2d_1d_2)n + (ab - ad_1 - bd_1 + d_1d_2)\]

 of the form $c_n = Ax^2 + Bx + C$, proving the claim.

 





\textbf{Solution 1 (faster)}

Use the same rationale as in solution 0.67; instead of using terms $1,2,3$, we use $-1,0,1$ and solve the $6$th term.

 $a-b+c=1440$

 $c=1716$

 $a+b+c=1848$

 Accordingly we will solve 

 $a=-72, b=204, c=1716$

 $36a+6b+c= \boxed{348}.$

 -maxamc

 





\textbf{Solution 2}

Setting one of the sequences as $a+nr_1$ and the other as $b+nr_2$, we can set up the following equalities

 $ab = 1440$

 $(a+r_1)(b+r_2)=1716$

 $(a+2r_1)(b+2r_2)=1848$

 We want to find $(a+7r_1)(b+7r_2)$

 Foiling out the two above, we have

 $ab + ar_2 + br_1 + r_1r_2 = 1716$ and $ab + 2ar_2 + 2br_1 + 4r_1r_2 = 1848$

 Plugging in $ab=1440$ and bringing the constant over yields

 $ar_2 + br_1 + r_1r_2 = 276$

 $ar_2 + br_1 + 2r_1r_2 = 204$

 Subtracting the two yields $r_1r_2 = -72$ and plugging that back in yields $ar_2 + br_1 = 348$

 Now we find

 $(a+7r_1)(b+7r_2) = ab + 7(ar_2 + br_1) + 49r_1r_2 = 1440 + 7(348) + 49(-72) = \boxed{348}$.

 





\textbf{Solution 3}

Let the first sequence be 

 $A={a+d_1, a + 2d_1, a + 3d_1, \cdots}$

 and the second be 

 $B={b+d_2, b + 2d_2, b + 3d_2, \cdots}$,

 with $(a+d_1)(b+d_2)=1440$. Now, note that the $n^{\text{th}}$ term of sequence $A$ is $a+d_1 n$ and the $n^{\text{th}}$ term of $B$ is $b + d_2  n$. Thus, the $n^{\text{th}}$ term of the given sequence is 

 $n^2(d_1 + d_2) + n(ad_2 + bd_1) + ab$,

 a quadratic in $n$. Now, letting the given sequence be $C$, we see that 

 $C_n - C_{n-1} = n^2(d_1 + d_2) + n(ad_2 + bd_1) + ab - (n-1)^2(d_1 + d_2) - (n-1)(ad_2 + bd_1) - ab = n(2d_1 + 2d_2) + ad_2 + bd_1 - d_1 - d_2$,

 a linear equation in $n$. Since $C_2 - C_1 = 276$ and $C_3 - C_2 = 132$, we can see that, in general, we have 

 $C_n - C_{n-1} = 420 - 144n$.

 Thus, we can easily find 

 $C_4 - C_3 = -12 \rightarrow C_4 = 1836$,

 $C_5 - C_4 = -156 \rightarrow C_5 = 1680$,

 $C_6 - C_5 = -300 \rightarrow C_6 = 1380$, 

 $C_7 - C_6 = -444 \rightarrow C_7 = 936$, and finally 

 $C_8 - C_7 = -588 \rightarrow \boxed{C_8 = 348}$.

 





\textbf{Solution 4(Tedious)}

Start by labeling the two sequences:

 Sequence 1:$a,a+d_1,a+2d_1,\dots a+(n-1)d_1$,

 Sequence 2:$b,b+d_2,b+2d_2,\dots b+(n-1)d_2$.

 Additionally, label the sequence given in the problem the function $f$, such that 

 $f(1)=1440,f(2)=1716,f(3)=1848$.

 Then, 

 $f(1)=ab,$

 $f(2)=(a+d_1)(b+d_2)=a+ad_2+d_1b+d_1d_2,$

 $f(3)=(a+2d_1)(b+2d_2)=ab+2ad_2+2bd_1+4d_1d_2,$

 and $f(8)=(a+7d_1)(b+7d_2)=ab+7ad_2+7bd_1+49d_1d_2$.

 In order to find $f(8)$ add $f(3)$ enough times to get the difference between the $d_1d_2$ and $ad_2+bd_1$ terms, then add $f(2)$ and $f(1)$ to get the other terms:

 $21f(3)=21ab+42ad_2+42bd_1+84d_1d_2$

 $21f(3)-35f(2)=21ab+42ad_2+42bd_1+84d_1d_2-35a-35ad_2-35d_1b-35d_1d_2=-14ab+7ad_2+7bd_1+49d_1d_2$

 $21f(3)-35f(2)+15f(1)=-14ab+7ad_2+7bd_1+49d_1d_2+15ab=ab+7ad_2+7bd_1+49d_1d_2$

 Now that the expression is in terms of the given values, insert values and solve:

 $21*1848-35*1716+15*1440=1848+20*132+(20*1716-35*1716+15*1716)+15*(-276)$ 

 $=1848+5*132+15(132-276)$

 $=1848+660+15(-144)$

 $=348$

 




	\item (\textit{AIME II 2003}) Consider the polynomials $P(x) = x^{6} - x^{5} - x^{3} - x^{2} - x$ and $Q(x) = x^{4} - x^{3} - x^{2} - 1.$ Given that $z_{1},z_{2},z_{3},$ and $z_{4}$ are the roots of $Q(x) = 0,$ find $P(z_{1}) + P(z_{2}) + P(z_{3}) + P(z_{4}).$


 
 \textbf{Answer: }006



\textbf{Solution 1}

When we use long division to divide $P(x)$ by $Q(x)$, the remainder is $x^2-x+1$.

 So, since $z_1$ is a root, $P(z_1)=(z_1)^2-z_1+1$.

 Now this also follows for all roots of $Q(x)$Now 
 \[P(z_2)+P(z_1)+P(z_3)+P(z_4)=z_1^2-z_1+1+z_2^2-z_2+1+z_3^2-z_3+1+z_4^2-z_4+1\]

 Now by \href{https://artofproblemsolving.com/wiki/index.php?title=Vieta\%27s\_formulas}{Vieta's} we know that $-z_4-z_3-z_2-z_1=-1$,so by \href{https://artofproblemsolving.com/wiki/index.php?title=Newton\%27s\_Sums}{Newton's Sums} we can find $z_1^2+z_2^2+z_3^2+z_4^2$

 $a_ns_2+a_{n-1}s_1+2a_{n-2}=0$

 $(1)(s_2)+(-1)(1)+2(-1)=0$

 $s_2-1-2=0$

 $s_2=3$

 So finally$P(z_2)+P(z_1)+P(z_3)+P(z_4)=3+4-1=\boxed{006}.$

 





\textbf{Solution 2 (Complete)}

By using long division to divide $P(x)$ by $Q(x)$, we find that $P(x) = Q(x)(x^2 + 1) + (x^2 - x + 1)$

 For \href{https://artofproblemsolving.com/wiki/index.php?title=Polynomial\_Remainder\_Theorem}{Polynomial Remainder Theorem},  
 \[P(z_2)+P(z_1)+P(z_3)+P(z_4)=z_1^2-z_1+1+z_2^2-z_2+1+z_3^2-z_3+1+z_4^2-z_4+1\] 

 
 \[\Rightarrow P(z_2)+P(z_1)+P(z_3)+P(z_4)=(z_1^2+z_2^2+z_3^2+z_4^2)-(z_1+z_2+z_3+z_4)+4\] 

 We know that for Vieta's Formula

 
 \[z_1+z_2+z_3+z_4=\frac{-b}{a}=\frac{-(-1)}{1}=1\]
 \[x_1x_2 + x_1x_3 + x_1x_4 + x_2x_3 + x_2x_4 + x_3x_4= \frac{c}{a} = \frac{-1}{1} = -1\]

 And, by Newton's Sums

 
 \[S=(z_1^2+z_2^2+z_3^2+z_4^2)= \sum_{i=1}^{n=4} z_i^2 = \left( \sum_{i=1}^{n=4} z_i \right)^2 - 2 \sum_{i=1}^{n=4}z_i z_j \text{ with } i<j\]

 So, 
 \[S= \sum_{i=1}^{n=4} z_i^2 = \left( \sum_{i=1}^{n=4} z_i \right)^2 - 2 \sum_{i=1}^{n=4}z_i z_j\]

 
 \[\Rightarrow S = (z_1+z_2+z_3+z_4)-2(x_1x_2 + x_1x_3 + x_1x_4 + x_2x_3 + x_2x_4 + x_3x_4)\]

 
 \[\Rightarrow S= 1 -2\times (-1) \Rightarrow S = 3\]

 Therefore

 
 \[\Rightarrow P(z_2)+P(z_1)+P(z_3)+P(z_4)=(z_1^2+z_2^2+z_3^2+z_4^2)-(z_1+z_2+z_3+z_4)+4\] 

 
 \[\Rightarrow P(z_2)+P(z_1)+P(z_3)+P(z_4) = (3)- 1 + 4 \therefore \boxed{ P(z_2)+P(z_1)+P(z_3)+P(z_4)=6}\]

 





\textbf{Solution 3}

Let $S_k=z_1^k+z_2^k+z_3^k+z_4^k$ then by \href{https://artofproblemsolving.com/wiki/index.php?title=Vieta\%27s\_Formula}{Vieta's Formula} we have
 \[S_{-1}=\frac{z_1z_2z_3+z_1z_3z_4+z_1z_2z_4+z_1z_2z_3}{z_1z_2z_3z_4}=0\]
 \[S_0=4\]
 \[S_1=1\]
 \[S_2=3\]By \href{https://artofproblemsolving.com/wiki/index.php?title=Newton\%27s\_Sums}{Newton's Sums} we have 
 \[a_4S_k+a_3S_{k-1}+a_2S_{k-2}+a_1S_{k-1}+a_0S_{k-4}=0\]

 Applying the formula couples of times yields $P(z_1)+P(z_2)+P(z_3)+P(z_4)=S_6-S_5-S_3-S_2-S_1=\boxed{6}$.

 





\textbf{Solution 4}

$P(x) = x^{2}Q(x)+x^{4}-x^{3}-x.$

 So we just have to find: $\sum_{n=1}^{4} z^{4}_n - \sum_{n=1}^{4} z^{3}_n - \sum_{n=1}^{4} z_n$.

 And by \href{https://artofproblemsolving.com/wiki/index.php?title=Newton\%27s\_Sums}{Newton's Sums} this computes to: $11-4-1 = \boxed{006}$.

 





\textbf{Solution 5}

If we scale $Q(x)$ by $x^2$, we get $x^6-x^5-x^4-x^2$. In order to get to $P(x)$, we add $x^4-x^3-x$. Therefore, our answer is $\sum_{n=1}^{4} z^4_n-z^3_n-z_n$. However, rearranging $Q(z_n) = 0$, makes our final answer $\sum_{n=1}^{4} z^2_n-z_n+1$. The sum of the squares of the roots is $1^2-2(-1) = 3$ and the sum of the roots is $1$. Adding 4 to our sum, we get $3-1+4 = \boxed{006}$.

 





\textbf{Solution 6}

Let $S_k$ = $z_1^k+z_2^k+z_3^k+z_4^k$

 By \href{https://artofproblemsolving.com/wiki/index.php?title=Newton\%27s\_Sums}{Newton's Sums},

 $S_1-1=0$

 $S_2-S_1-2=0$

 $S_3-S_2-S_1=0$

 $S_4-S_3-S_2-4=0$

 $S_5-S_4-S_3-S_1=0$

 $S_6-S_5-S_4-S_2=0$

 Solving for $S_1,S_2,S_3,S_4,S_5,S_6$, we get $S_1=1, S_2=3, S_3=4, S_4=11, S_5=16, S_6=30$

 $P(z_1)+P(z_2)+P(z_3)+P(z_4)=S_6-S_5-S_3-S_2-S_1=\boxed{006}$

 





\textbf{Solution 7 (if you don't know how to divide polynomials directly)}

$P(x)$ and $Q(x)$  look very similar, so let's try subtracting $Q(x)$ multiplied with something from $P(x)$ since we know $Q(z_n) = 0$. To make the degrees the same, let's first multiply $Q(x)$ with $x^2$: 

 $Q(x) \cdot x^2 = x^6 - x^5 - x^4 - x^2$. Subtracting this from $P(x)$ gives $x^4 - x^3 - x$. The degrees of this polynomial and $Q(x)$ are the same, so let's subtract $Q(x)$: $x^4 - x^3 - x - Q(x)$ = $x^4 - x^3 - x - x^4 + x^3 + x^2 + 1 = x^2 - x + 1.$

 So $P(x) - (x^2 + 1) \cdot Q(x) = x^2 - x + 1$ (this is also the remainder when $P(x)$ is divided by $Q(x)$). This tells us that our answer is the sum of the squares of the roots of $Q(x)$ minus the sum of the roots of $Q(x)$ plus $1 \cdot 4$.  

 By \href{https://artofproblemsolving.com/wiki/index.php?title=Newton\%27s\_Sums}{Newton's Sums} the sum of the squares of the roots is 3, and by \href{https://artofproblemsolving.com/wiki/index.php?title=Vieta\%27s\_Formulas}{Vieta's Formulas} the sum of the roots is 1. So our answer is $3 - 1  + 4 = \boxed{006}$. 

 
[rule]

 Nice!-sleepypuppy

 




	\item (\textit{AIME II 2003}) Two positive integers differ by $60.$ The sum of their square roots is the square root of an integer that is not a perfect square. What is the maximum possible sum of the two integers?


 
 \textbf{Answer: }156



\textbf{Solution 1}

Call the two integers $b$ and $b+60$, so we have $\sqrt{b}+\sqrt{b+60}=\sqrt{c}$. Square both sides to get $2b+60+2\sqrt{b^2+60b}=c$. Thus, $b^2+60b$ must be a square, so we have $b^2+60b=n^2$, and $(b+n+30)(b-n+30)=900$. The sum of these two factors is $2b+60$, so they must both be even. To maximize $b$, we want to maximixe $b+n+30$, so we let it equal $450$ and the other factor $2$, but solving gives $b=196$, which by a quick check yields perfect square $c=900$, so we have to keep going. In order to keep both factors even, we let the larger one equal $150$ and the other $6$, which gives $b=48$. This checks, so the solution is $48+108=\boxed{156}$.

 




	\item (\textit{AIME II 2003}) Triangle $ABC$ is a right triangle with $AC = 7,$ $BC = 24,$ and right angle at $C.$ Point $M$ is the midpoint of $AB,$ and $D$ is on the same side of line $AB$ as $C$ so that $AD = BD = 15.$ Given that the area of triangle $CDM$ may be expressed as $\frac {m\sqrt {n}}{p},$ where $m,$ $n,$ and $p$ are positive integers, $m$ and $p$ are relatively prime, and $n$ is not divisible by the square of any prime, find $m + n + p.$


 
 \textbf{Answer: }578



\textbf{Solution 1}

We use the \href{https://artofproblemsolving.com/wiki/index.php?title=Pythagorean\_Theorem}{Pythagorean Theorem} on $ABC$ to determine that $AB=25.$

 Let $N$ be the orthogonal projection from $C$ to $AB.$ Thus, $[CDM]=\frac{(DM)(MN)} {2}$, $MN=BN-BM$, and $[ABC]=\frac{24 \cdot 7} {2} =\frac{25 \cdot (CN)} {2}.$

 From the third equation, we get $CN=\frac{168} {25}.$

 By the \href{https://artofproblemsolving.com/wiki/index.php?title=Pythagorean\_Theorem}{Pythagorean Theorem} in $\Delta BCN,$ we have 

 $BN=\sqrt{\left(\frac{24 \cdot 25} {25}\right)^2-\left(\frac{24 \cdot 7} {25}\right)^2}=\frac{24} {25}\sqrt{25^2-7^2}=\frac{576} {25}.$ 

 Thus, $MN=\frac{576} {25}-\frac{25} {2}=\frac{527} {50}.$ 

 In $\Delta ADM$, we use the \href{https://artofproblemsolving.com/wiki/index.php?title=Pythagorean\_Theorem}{Pythagorean Theorem} to get $DM=\sqrt{15^2-\left(\frac{25} {2}\right)^2}=\frac{5} {2} \sqrt{11}.$

 Thus, $[CDM]=\frac{527 \cdot 5\sqrt{11}} {50 \cdot 2 \cdot 2}= \frac{527\sqrt{11}} {40}.$

 Hence, the answer is $527+11+40=\boxed{578}.$

 





\textbf{Solution 2}

By the \href{https://artofproblemsolving.com/wiki/index.php?title=Pythagorean\_Theorem}{Pythagorean Theorem} in $\Delta AMD$, we get $DM=\frac{5\sqrt{11}} {2}$. Since $ABC$ is a right triangle, $M$ is the circumcenter and thus, $CM=\frac{25} {2}$. We let $\angle CMD=\theta$. By the \href{https://artofproblemsolving.com/wiki/index.php?title=Law\_of\_Cosines}{Law of Cosines},

 $24^2 = 2 \cdot (12.5)^2-2 \cdot (12.5)^2 * \cos (90+\theta).$

 It follows that $\sin \theta = \frac{527} {625}$. Thus,$[CMD]=\frac{1} {2} (12.5) \left(\frac{5\sqrt{11}} {2}\right)\left(\frac{527} {625}\right)=\frac{527\sqrt{11}} {40}$.

 





\textbf{Solution 3}

Suppose $ABC$ is plotted on the \href{https://artofproblemsolving.com/wiki/index.php?title=Cartesian\_plane}{cartesian plane} with $C$ at $(0,0)$, $A$ at $(0,7)$, and $B$ at $(24,0)$.Then $M$ is at $(12,3.5)$. Since $\Delta ABD$ is isosceles, $MD$ is perpendicular to $AM$, and since $AM=12.5$ and $AD=15, MD=2.5\sqrt{11}$.The slope of $AM$ is $-\frac{7}{24}$ so the slope of $MD$ is $\frac{24}{7}$.Draw a vertical line through $M$ and a horizontal line through $D$. Suppose these two lines meet at $X$. then $MX=\frac{24}{7}DX$ so $MD=\frac{25}{7}DX=\frac{25}{24}MD$ by the pythagorean theorem.So $MX=2.4\sqrt{11}$ and $DX=.7\sqrt{11}$ so the coordinates of D are $(12-.7\sqrt{11},\, 3.5-2.4\sqrt{11})$.Since we know the coordinates of each of the vertices of $\Delta CMD$, we can apply the \href{https://artofproblemsolving.com/wiki/index.php?title=Shoelace\_Theorem}{Shoelace Theorem} to find the area of $\Delta CMD, \frac{527 \sqrt{11}}{40}$.

 





\textbf{Solution 4}

Let $E$ be the intersection of lines $BC$ and $DM$.  Since triangles $\Delta CME$ and $\Delta CMD$ share a side and height, the area of $\Delta CDM$ is equal to $\frac{DM}{EM}$ times the area of $\Delta CME$.By AA similarity, $\Delta EMB$ is similar to $\Delta ACB$, $\frac{EM}{AC}=\frac{MB}{CB}$. Solving yields $EM=\frac{175}{48}$. Using the same method but for $EB$ yields $EB=\frac{625}{48}$. As in previous solutions, by the \href{https://artofproblemsolving.com/wiki/index.php?title=Pythagorean\_Theorem}{Pythagorean Theorem}, $DM=\frac{5\sqrt{11}}{2}$. So, $\frac{DM}{EM}=\frac{24\sqrt{11}}{35}$.Now, since we know both $CB$ and $EB$, we can find that $CE=\frac{527}{48}$. The height of $\Delta CME$ is the length from point $M$ to $CB$. Since $M$ is the midpoint of $AB$, the height is just $\frac{1}{2}\cdot7=\frac{7}{2}$. Using this, we can find that the area of $\Delta CMD$ is $\frac{1}{2}\cdot\frac{7}{2}\cdot\frac{527}{48}\cdot\frac{24\sqrt{11}}{35}=\frac{527\sqrt{11}}{40}$, giving our answer of $\boxed{578}$.

 Solution by someonenumber011.

 





\textbf{Solution 5 [braindead simple + asy diagram]}


                    \begin{figure}[H]
                        \centering
                        \includegraphics[width=0.5\textwidth]{Solution/AIME_2003_II_11_Solution_Figure_1.png}
                    \end{figure}
                    First, notice that 
 \[BM = AM = \frac{25}{2}.\] Further, notice that 
 \[[ABC] = \frac{7*24}{2}=84. \frac{[\triangle{AMC}]}{\triangle{ABC}}=\frac{AM}{AB}=\frac{1}{2}\]
 \[\therefore [\triangle{AMC}]=42\]Then, also notice that 
 \[[AMDC] =[\triangle{AMD}]+[\triangle{ACD}]=[\triangle{AMC}]+[\triangle{CMD}]=42+[\triangle{CMD}]\]So we can say:
 \[[\triangle{CMD}]=[\triangle{AMD}]+[\triangle{ACD}]-42\]Then, notice that $[\triangle{AMD}] = AM * MD /2,$ since D lies on the perpindicular bisector of $\overline{AB}$ (because $AD=BD=15$) and therefore $\overline{AM} \perp \overline{MD}$.
 \[MD  = \sqrt{15^2-\frac{25^2}{4}}= \frac{5\sqrt{11}}{2}\implies [\triangle{AMD}]=\frac{25}{2}*\frac{5\sqrt{11}}{2}*\frac{1}{2}=\frac{125\sqrt{11}}{8}\]Now, take $\alpha = \angle{CAB},\beta = \angle{DAM},\gamma = \beta-\alpha$. Then, 
 \[\sin(\beta) = \frac{\sqrt{11}}{6},\cos(\beta) =\frac{5}{6},\sin(\alpha) =24/25,\cos(\alpha)=7/25\]So, 
 \[\sin(\gamma)=\sin(\alpha-\beta)=\frac{24}{25}*\frac{5}{6} - \frac{7}{25}*\frac{\sqrt{11}}{6}=\frac{120-7\sqrt{11}}{150}\]
 \[[\triangle{ACD}]=\frac{7*15*\sin(\gamma)}{2}=\frac{7}{2}*\frac{120-7\sqrt{11}}{10}=42-\frac{49\sqrt{11}}{20}\]and thus:
 \[[\triangle{CMD}]=42-\frac{49\sqrt{11}}{20}+\frac{125\sqrt{11}}{8}-42=\frac{625\sqrt{11}-98\sqrt{11}}{40}=\frac{527\sqrt{11}}{40}\implies m+n+p = 527+11+40=\boxed{\textbf{578}}\]

 




	\item (\textit{AIME II 2003}) The members of a distinguished committee were choosing a president, and each member gave one vote to one of the $27$ candidates. For each candidate, the exact percentage of votes the candidate got was smaller by at least $1$ than the number of votes for that candidate. What is the smallest possible number of members of the committee?


 
 \textbf{Answer: }134



\textbf{Solution}

Let $v_i$ be the number of votes candidate $i$ received, and let $s=v_1+\cdots+v_{27}$ be the total number of votes cast. Our goal is to determine the smallest possible $s$.

 Candidate $i$ got $\frac{v_i}s$ of the votes, hence the percentage of votes they received is $\frac{100v_i}s$. The condition in the problem statement says that $\forall i: \frac{100v_i}s + 1 \leq v_i$.  ($\forall$ means "for all", so this means "For all $i$, $\frac{100v_i}s + 1 \leq v_i$ is true")

 Obviously, if some $v_i$ would be $0$ or $1$, the condition would be false. Thus $\forall i: v_i\geq 2$. We can then rewrite the above inequality as $\forall i: s\geq\frac{100v_i}{v_i-1}$. 

 If for some $i$ we have $v_i=2$, then from the inequality we just derived we would have $s\geq 200$.

 If for some $i$ we have $v_i=3$, then $s\geq 150$.

 And if for some $i$ we have $v_i=4$, then $s\geq \frac{400}3 = 133\frac13$, and hence $s\geq 134$.

 Is it possible to have $s<134$? We just proved that to have such $s$, all $v_i$ have to be at least $5$. But then $s=v_1+\cdots+v_{27}\geq 27\cdot 5 = 135$, which is a contradiction. Hence the smallest possible $s$ is at least $134$.

 Now consider a situation where $26$ candidates got $5$ votes each, and one candidate got $4$ votes. In this situation, the total number of votes is exactly $134$, and for each candidate the above inequality is satisfied. Hence the minimum number of committee members is $s=\boxed{134}$.

 Note: Each of the $26$ candidates received $\simeq 3.63\%$ votes, and the last candidate received $\simeq 2.985\%$ votes.

 





\textbf{Solution 2}

Let there be $N$ members of the committee.Suppose candidate $n$ gets $a_n$ votes.Then $a_n$ as a percentage out of $N$ is $100\frac{a_n}{N}$. Setting up the inequality $a_n \geq 1 + 100\frac{a_n}{N}$ and simplifying,$a_n \geq \lceil(\frac{N}{N - 100})\rceil$ (the ceiling function is there because $a_n$ is an integer.Note that if we set all $a_i$ equal to $\lceil(\frac{N}{N - 100})\rceil$ we have $N \geq 27\lceil(\frac{N}{N - 100})\rceil$. Starting from $101$, try greater values of $N$ to find which value is large enough. It helps to try $N$ with multiples of $5$ or $10$ for ease of computation.Then when $N=125$, the RHS equals $27(5)=135$, which is pretty close. As soon as the ceiling part equals $4$, RHS equals $108\le N$. Solving $\frac{N}{N - 100}\leq4$ yields $N = 134$ as the least such number that satisfies this inequality. Now we must show that we can find suitable $a_i$. We can let 26 of them equal to $5$ and one of them equal to $4$ or 26 equal $4$ and one equal $30$. Therefore, $N = \boxed{134}$ is the answer.- whatRthose

 





\textbf{Solution 3}

Let $n$ be the total number of people in the committee, and $a_i$ be the number of votes candidate $i$ gets where $1 \leq i \leq 27$.  The problem tells us that 
 \[\frac{100a_i}{n} \leq a_i - 1 \implies 100a_i \leq na_i - n \implies a_i \geq \frac{n}{n-100}.\]Therefore, 
 \[\sum^{27}_{i=1} a_i = n \geq \sum^{27}_{i=1} \frac{n}{n-100} = \frac{27n}{n-100},\]and so $n(n-127) \geq 0 \implies n \geq 127$.  Trying $n = 127$, we get that 
 \[a_i \geq \frac{127}{27} \approx 4.7 \implies a_i \geq 5 \implies \sum^{27}_{a_i} a_i \geq 5 \cdot 27 = 135 \geq 127,\]a contradiction.  Bashing out a few more, we find that $\boxed{n = 134}$ works perfectly fine.

 




	\item (\textit{AIME II 2003}) A bug starts at a vertex of an equilateral triangle. On each move, it randomly selects one of the two vertices where it is not currently located, and crawls along a side of the triangle to that vertex. Given that the probability that the bug moves to its starting vertex on its tenth move is $m/n,$ where $m$ and $n$ are relatively prime positive integers, find $m + n.$


 
 \textbf{Answer: }683



\textbf{Solution 1A (Easiest)}

Let $P_n$ represent the probability that the bug is at its starting vertex after $n$ moves. If the bug is on its starting vertex after $n$ moves, then it must be not on its starting vertex after $n-1$ moves. At this point it has $\frac{1}{2}$ chance of reaching the starting vertex in the next move. Thus $P_n=\frac{1}{2}(1-P_{n-1})$. $P_0=1$, so now we can build it up:

 $P_1=0$,$P_2=\frac{1}{2}$,$P_3=\frac{1}{4}$,$P_4=\frac{3}{8}$,$P_5=\frac{5}{16}$,$P_6=\frac{11}{32}$,$P_7=\frac{21}{64}$,$P_8=\frac{43}{128}$,$P_9=\frac{85}{256}$,$P_{10}=\frac{171}{512}$,

 Thus the answer is $171+512=$$\boxed{683}$

 





\textbf{Solution 1B (very detailed with explanations, ~grogg007)}

Let $Q_{N}$ be the probability of starting at one point and ending at another for any pair of points on the triangle after $N$ moves and $P_{N}$ be the probability of ending at the same point you started on for any point on the triangle after $N$ moves. 

 You can either end on two different points you started from or end at the same point you started from, so $Q_{N} + Q_{N} + P_{N} = 1.$ Now, let's write $P_{10}$ as a sequence of moves. Say we start at point $A$ of the triangle and the vertices are labeled $A, B,$ and $C.$ Let each arrow represent a move to a vertex:

 
 \[A \rightarrow \_\_\_\_ \rightarrow \_\_\_\_ \rightarrow \_\_\_\_ \dots \rightarrow \_\_\_\_ \rightarrow A\]

 We have two main cases. Case 1: both the second slot and ninth slot are the same vertex (WLOG say it's vertex $B$): 
 \[A \rightarrow B \rightarrow \_\_\_\_ \rightarrow \_\_\_\_ \dots \rightarrow B \rightarrow A\]

 Then this is simply $P_{8}$ to get from B on the second move to B again on the ninth move multiplied with the probability of going from $B$ to $A$ on the tenth move, which is $1/2.$ So this case gives us $P_{8}/2$ in our probability. The probability of going from $A$ to $B$ on the first move is $1/2,$ but since $B$ can also be $C,$ we multiply this by $2$ to just get $P_{8}/2 \cdot 1/2 \cdot 2 = P_{8}/2$.

 Case 2: the second and ninth slots are two different vertices (WLOG say it's vertex $B$ in the second slot and vertex $C$ in the ninth slot):
 \[A \rightarrow B \rightarrow \_\_\_\_ \rightarrow \_\_\_\_ \dots \rightarrow C \rightarrow A\]

 Then this is just $Q_{8}$ to get from B on the second move to C on the ninth move, along with the probability of going form $C$ to $A$ on the tenth move, which is $1/2.$ The probability of going from $A$ to $B,$ like in the first case, is $1/2,$ but we could also have $C$ in the second slot and $B$ in the ninth slot, so like case $1$ we multiply this by $2$ to just get $1/2 \cdot 2 \cdot Q_{8} \cdot 1/2 = Q_{8}/2.$ 

 Together, we have $P_{10} = \frac{P_{8} + Q_{8}}{2}.$ We can repeat this process for all $N$ and generalize by saying $P_{N} = \frac{P_{N-2} + Q_{N-2}}{2}.$

 Now, since $Q_{N} + Q_{N} + P_{N} = 2 \cdot Q_{N} + P_{N} = 1,$ we know that $2 \cdot Q_{N-2} + P_{N-2} = 1.$ Solving for $Q_{N-2}$ we get $Q_{N-2} = \frac{1 - P_{N-2}}{2}.$ Plugging this back into our equation for $P_{N}$ gives us $P_{N} = \frac{P_{N-2} + 1}{4}.$ From here, proceed with recursion to find $P_{10}$ (which is the answer to the problem). We need a value of $P_{N}$ to start the process. We will find $P_{2}.$ Since you would have to go back and forth between a point with two moves to end back where you started with probability $\frac{1}{4},$ but you could also select any of the two points you aren't already at, the probability is $\frac{1}{4} \cdot 2 = \frac{1}{2}.$ So $P_{2} = \frac{1}{2}.$ We work up from here: 
 \[P_{4} = \frac{\frac{1}{2} + 1}{4} = \frac{3}{8}\]
 \[P_{6} = \frac{\frac{3}{8} + 1}{4} = \frac{11}{32}\] 
 \[P_{8} = \frac{\frac{11}{32} + 1}{4} = \frac{43}{128}\] 
 \[P_{10} = \frac{\frac{43}{128} + 1}{4} = \boxed{\frac{171}{512}}.\]

 $171 + 512 = \boxed{683}.$

 





\textbf{Solution 2 (also easiest)}

Consider there to be a clockwise and a counterclockwise direction around the triangle.  Then, in order for the ant to return to the original vertex, the net number of clockwise steps must be a multiple of 3, i.e., $\#CW - \#CCW \equiv 0 \pmod{3}$.  Since $\#CW + \#CCW = 10$, it is only possible that $(\#CW,\, \#CCW) = (5,5), (8,2), (2,8)$.

 In the first case, we pick $5$ out of the ant's $10$ steps to be clockwise, for a total of ${10 \choose 5}$ paths.  In the second case, we choose $8$ of the steps to be clockwise, and in the third case we choose $2$ to be clockwise.  Hence the total number of ways to return to the original vertex is ${10 \choose 5} + {10 \choose 8} + {10 \choose 2} = 252 + 45 + 45 = 342$.  Since the ant has two possible steps at each point, there are $2^{10}$ routes the ant can take, and the probability we seek is $\frac{342}{2^{10}} = \frac{171}{512}$, and the answer is $\boxed{683}$.

 





\textbf{Solution 3}

Label the vertices of the triangle $A,B,C$ with the ant starting at $A$. We will make a table of the number of ways to get to $A,B,C$ in $n$ moves $n\leq10$. The values of the table are calculated from the fact that the number of ways from a vertex say $A$ in $n$ steps equals the number of ways to get to $B$ in $n-1$ steps plus the number of ways to get to $C$ in $n-1$ steps.

 


 

 \[\begin{array}{|l|ccc|} \multicolumn{4}{c}{\text{Table}}\\\hline \text{Step}&A&B&C \\\hline 1 &0 &1 &1 \\ 2 &2 &1 &1 \\ 3 &2 &3 &3\\ \vdots &\vdots&\vdots&\vdots \\ 10 &342 &341 &341 \\\hline \end{array}\]Therefore, our answer is $512+171=\boxed{683}.$

 


 Notice the pattern that there are $\left\lceil\frac{2^n}{3}\right\rceil$ way to get to $A$ for even $n$ moves. Thus, there are $\left\lceil\frac{2^{10}}{3}\right\rceil=342$ ways.

 





\textbf{Solution 4}

Notice that this problem can be converted into a Markov Chain transition matrix.

 The transition matrix is { {0,1,1}, {1,0,1} , {1,1,0} } * (1/2) . Then use the exponentiation method of squaring ( A*A---(A\textasciicircum{}2)*(A\textasciicircum{}2)---(A\textasciicircum{}4*A\textasciicircum{}4)--(A\textasciicircum{}8*A\textasciicircum{}2) to get the transition value of 342. Divide by 2\textasciicircum{}10 for the probability, reduce fractions, for the result of 171+512 = 683.

 





\textbf{Solution 4.1 (solution 4 but rigorized and generalized)}

As a note, do NOT do this on the exam as it will eat up your time, but feel free to experiment around with this if you have a good enough understanding of linear algebra. This writeup will be extremely lengthy because I am assuming just very basic concepts of linear algebra. The concepts here extend to higher levels of mathematics, so feel free to explore more in depth so that you can end up solving almost any variation of this classic problem.

 \textbf{4.1.1: Conceptual Setup} There are a possible of 3 states for this problem, so we can model the problem as a stochastic process. The resulting process has a transition matrix of:

 
 \[\hat{T} = \begin{bmatrix} 0 & 0.5 & 0.5\\ 0.5 & 0 & 0.5\\ 0.5 & 0.5 & 0 \end{bmatrix}\]

 We aim to diagonalize this transition matrix to make it easier to exponentiate by converting it into what's known as it's Jordan Canonical Form. 

 \textbf{4.1.2: Eigenvalue/vector bash} In order to do this, we must extract the eigenvalues and eigenvectors of the matrix. The eigenpolynomial for this matrix is obtained by calculating this matrix's determinant with $0-\lambda$ about it's main diagonal like so:
 \[\hat{T}_{\lambda} = \begin{bmatrix} 0-\lambda & 0.5 & 0.5\\ 0.5 & 0-\lambda & 0.5\\ 0.5 & 0.5 & 0-\lambda \end{bmatrix}\]

 We have the matrix's eigenpolynomial to be $\lambda^3 - \frac{3\lambda}{4} + \frac{1}{4}$, and extracting eigenvalues by setting the polynomial equal to $0$, we have 2 eigenvalues: $\lambda_1 = 1$ of multiplicity 1, and $\lambda_2 = -\frac{1}{2}$ of multiplicity 2. To extract the eigenvectors, we must assess the kernel of this matrix (also known as the null space), or the linear subspace of the domain of $\hat{T}$ where everything gets mapped to the null vector. 

 We first do this for $\lambda_1$. Taking $-\lambda_1$ across the diagonals to get $\hat{T}_{\lambda_1} = \begin{bmatrix} -1 & 0.5 & 0.5\\ 0.5 & -1 & 0.5\\ 0.5 & 0.5 & -1 \end{bmatrix}$, we first reduce it to reduced row echelon form, which is $\begin{bmatrix} 1 & 0 & -1\\ 0 & 1 & -1\\ 0 & 0 & 0 \end{bmatrix}$. From here, we compute the kernel by setting  
 \[\begin{bmatrix} 1 & 0 & -1\\ 0 & 1 & -1\\ 0 & 0 & 0 \end{bmatrix} \begin{bmatrix} x_1\\ x_2\\ x_3  \end{bmatrix} = 0\]. So if we take our free variable $x_3 = 0 = t$, then that means that in the same fashion, $x_1 - x_1 = x_2 - x_2 = 0 = t$, so hence, the kernel of $\hat{T}_{\lambda} = \begin{bmatrix} t\\ t\\ t \end{bmatrix}$, or more simply, $t\begin{bmatrix} 1\\ 1\\ 1 \end{bmatrix}$. $\begin{bmatrix} 1\\ 1\\ 1 \end{bmatrix}$ is the eigenvector corresponding to $\lambda_1$. We do the same computations for our second unique eigenvalue, but I will save the computation to you. There are actually 2 eigenvectors for $\lambda_2$, because the reduced row echelon form for $\hat{T}_{\lambda_2}$ has 2 free variables instead of 1, so our eigenvectors for $\lambda_2$ are $\begin{bmatrix} -1\\ 1\\ 0 \end{bmatrix}, \begin{bmatrix} -1\\ 0\\ 1 \end{bmatrix}$. We are now ready to begin finding the Jordan canonical form

 \textbf{4.1.3 Jordan Canonical Form} In linear algebra, the JCF (which also goes by the name of Jordan Normal Form) is an upper triangular matrix representing a linear transformation over a finite-dimensional complex vector space. Any square matrix that has a Jordan Canonical Form has its field of coefficients extended into a field containing all it's eigenvalues. You can find more information about them on google, as well as exactly how to find them but for now let's get on with the problem. I will skip the computation in this step, largley because this writeup is already gargantuan for a simple AIME problem, and because there are countless resources explaining how to do so.

 We aim to decompose $\hat{T}$ into the form $\hat{T} = SJS^{-1}$, where $S$ is a matrix whose columns consist of the eigenvectors of $\hat{T}$, $J$ is the Jordan matrix, and $S^{-1}$ is, well, the inverse of $S$. We have 1 eigenvalue of multiplicity 1, and 1 of multiplicity 2, so based on this info, we set our eigenvalues along the diagonals like so.We have:
 \[J = \begin{bmatrix} -\frac{1}{2} & 0 & 0\\ 0 & -\frac{1}{2} & 0\\ 0 & 0 & 1 \end{bmatrix}, S = \begin{bmatrix} -1 & -1 & 1\\ 0 & 1 & 1\\ 1 & 0 & 1 \end{bmatrix}, S^{-1} = \frac{1}{3}\begin{bmatrix} -1 & -1 & 2\\ -1 & 2 & -1\\ 1 & 1 & 1 \end{bmatrix}\]and so:
 \[\hat{T} = \begin{bmatrix} 0 & 0.5 & 0.5\\ 0.5 & 0 & 0.5\\ 0.5 & 0.5 & 0 \end{bmatrix} = SJS^{-1} = \begin{bmatrix} -1 & -1 & 1\\ 0 & 1 & 1\\ 1 & 0 & 1 \end{bmatrix}\begin{bmatrix} -\frac{1}{2} & 0 & 0\\ 0 & -\frac{1}{2} & 0\\ 0 & 0 & 1 \end{bmatrix}\begin{bmatrix} -\frac{1}{3} & -\frac{1}{3} & \frac{2}{3}\\ -\frac{1}{3} & \frac{2}{3} & -\frac{1}{3}\\ \frac{1}{3} & \frac{1}{3} & \frac{1}{3} \end{bmatrix}\]Now that we have converted to Jordan Canonical Form, it is extremely easy to compute $\hat{T}^n$.

 \textbf{4.1.4: Using the JCF to calculate the transition matrix to the power of any n, large or small} It is an important fact that for any matrix $K$ with Jordan decomposition $SJ_kS^{-1}$, we have that $K^n = S(J_k)^nS^{-1}$Using this fact, we aim to find the general solution for the problem. $J^n = \begin{bmatrix} \left(-\frac{1}{2}\right)^n & 0 & 0\\ 0 & \left(-\frac{1}{2}\right)^n & 0\\ 0 & 0 & 1 \end{bmatrix}$, and using the laws of matrix multiplication, 
 \[SJ^n = \begin{bmatrix} (-1)^{n+1}\left(\frac{1}{2}\right)^n & (-1)^{n+1}\left(\frac{1}{2}\right)^n & 1\\ 0 & \left(-\frac{1}{2}\right)^n & 1\\ \left(-\frac{1}{2}\right)^n & 0 & 1 \end{bmatrix}\]So finally, our final, generalized transition matrix after any number of steps $n$ is:
 \[\hat{T}^n = \begin{bmatrix} \frac{1}{3} - \frac{1}{3}(-1)^{1+n}(2)^{1-n} & \frac{1}{3} + \frac{1}{3}(-1)^{n+1}(2)^{1-(1+n)} & \frac{1}{3} + \frac{1}{3}(-1)^{n+1}(2)^{1-(1+n)}\\ \frac{1}{3} + \frac{1}{3}(-1)^{n+1}(2)^{1-(1+n)} & \frac{1}{3} - \frac{1}{3}(-1)^{1+n}(2)^{1-n} & \frac{1}{3} + \frac{1}{3}(-1)^{n+1}(2)^{1-(1+n)}\\ \frac{1}{3} + \frac{1}{3}(-1)^{n+1}(2)^{1-(1+n)} & \frac{1}{3} + \frac{1}{3}(-1)^{n+1}(2)^{1-(1+n)} & \frac{1}{3} - \frac{1}{3}(-1)^{1+n}(2)^{1-n} \end{bmatrix}\]

 \textbf{4.1.5 just plugging in for this problem lol} For the sake of this problem, we seek the top left element, which is $\frac{1}{3} - \frac{1}{3}(-1)^{1+n}(2)^{1-n}$. Substituting $n=10$ readily gives the probability of the bug reaching it's starting position within 10 moves to be $\frac{171}{512} \implies m+n = \boxed{683}$

 And we have also derived formulae for the bug reaching any state within $n$ moves as a byproduct!

 





\textbf{Solution 5 (guess & check)}

This method does not rigorously get the answer, but it works. As the bug makes more and more moves, the probability of it going back to the origin approaches closer and closer to 1/3. Therefore, after 10 moves, the probability gets close to $341.33/1024$. We can either round up or down. If we round down, we see $341/1024$ cannot be reduced any further and because the only answers on the AIME are below 1000, this cannot be the right answer. However, if we round up, $342/1024$ can be reduced to $171/512$ where the sum 171+512= $\boxed{683}$ is an accepted answer.

 





\textbf{Solution 6 (generating functions)}

The generating function for this is $(x+x^2)$ since an ant on any vertex of the equilateral triangle can go $120$ degrees or $240$ degrees to a side and simplifying $(x^{120}+x^{240})$ gets you $(x+x^2)$. Since $360$ degrees brings you back to the original vertex then we must find the sum of the coefficients that share a variable with a power divisible by $3$.

 Since we take this rotation $10$ times, our function becomes $(x+x^2)^{10}$ which is the same as $x^{10}(x+1)^{10}$. This completely simplified is $x(x+1)^{10}$ and since your maximum power is $11$, we only have to find the coefficients for $3$, $6$, and $9$ ($0$ would apply here but the $x$ is the lowest power there is).

 For $x^9$, the coefficient is ${10 \choose 2}$ , and the same goes for $x^3$. For $x^6$, the coefficient is ${10 \choose 5}$ and the final sum for the numerator is $2*{10 \choose 2} + {10 \choose 5}$ . The total sum is $90+252=342$ and for the denominator, it was simply $2^{10}$ and this simplified was $171/512$. Therefore the sum is $683$.

 





\textbf{Solution 7 (trees)}

Start of with any vertex, say $A$. Denote $a_n$ the number of paths ending where it started. Then notice that for a path to end in the vertex you started on, you must have for the $(n-1)$ case that of the $2^{n-1}$ total paths, we must take away the paths that end in the $(n-1)$-st term where it started off. (Because imagine on the $(n-1)$ move that it landed on vertex $A$. Then if we wanted to find the number of length $n$ paths that end at $A$, we would be able to find that path, because on the $(n-1)$-st move it landed at $A$. You can't go from $A$ to $A$). Hence we have the recursion $a_n=2^{n-1}-a_{n-1}$, with $a_3 = 2$. Now reiterating gives us $a_{10} = 342$, so that the probability is $\frac{342}{2^{10}} = \frac{171}{512}$. So we have $171 + 512 = \boxed{683}$.~th1nq3r

 (Note: One might be confused because you might think maybe "But we only did it for case $A$. Now for $B$ and $C$. Oh wait they are symmetric. So then if this is the correct answer, why am I wrong, or what happened to that factor of $3$?" Well truth be told, I skipped a somewhat major step. Notice that the total amount $2^n$ only comes from if we START OFF at vertex $A$. So we really need to multiply $2^n$ by $3$ to get $3(2^n)$ as the TRUE total amount. However this factor of $3$ doesn't matter because now that we do case work, let $a_n$ once again denote the number of paths starting at vertex $A$. Then abusing symmetry, we have that the number of paths ending where it started as $3(a_n)=3(2^{n-1}-a_{n-1})$. So now when we take the probability, we cancel the factor of $3$, and obtain the same answer as if it were for all three vertices).

 \textbf{Video Solution by Sal Khan} 




	\item (\textit{AIME II 2003}) Let $A = (0,0)$ and $B = (b,2)$ be points on the coordinate plane. Let $ABCDEF$ be a convex equilateral hexagon such that $\angle FAB = 120^\circ,$ $\overline{AB}\parallel \overline{DE},$ $\overline{BC}\parallel \overline{EF,}$ $\overline{CD}\parallel \overline{FA},$ and the y-coordinates of its vertices are distinct elements of the set $\{0,2,4,6,8,10\}.$ The area of the hexagon can be written in the form $m\sqrt {n},$ where $m$ and $n$ are positive integers and n is not divisible by the square of any prime. Find $m + n.$


 
 \textbf{Answer: }051



\textbf{Solution 1}

The y-coordinate of $F$ must be $4$.  All other cases yield non-convex and/or degenerate hexagons, which violate the problem statement.

 Letting $F = (f,4)$, and knowing that $\angle FAB = 120^\circ$, we can use rewrite $F$ using complex numbers:$f + 4 i = (b + 2 i)\left(e^{i(2 \pi / 3)}\right) = (b + 2 i)\left(-\frac{1}{2} + \frac{\sqrt{3}}{2} i\right) = -\frac{b}{2}-\sqrt{3}+\left(\frac{b\sqrt{3}}{2}-1\right)i$.  We solve for $b$ and $f$ and find that $F = \left(-\frac{8}{\sqrt{3}}, 4\right)$ and that $B = \left(\frac{10}{\sqrt{3}}, 2\right)$.

 The area of the hexagon can then be found as the sum of the areas of two congruent triangles ($EFA$ and $BCD$, with height $8$ and base $\frac{8}{\sqrt{3}}$) and a parallelogram ($ABDE$, with height $8$ and base $\frac{10}{\sqrt{3}}$).

 $A = 2 \times \frac{1}{2} \times 8 \times \frac{8}{\sqrt{3}} + 8 \times \frac{10}{\sqrt{3}} = \frac{144}{\sqrt{3}} = 48\sqrt{3}$.

 Thus, $m+n = \boxed{051}$.

 





\textbf{Solution 2}


                    \begin{figure}[H]
                        \centering
                        \includegraphics[width=0.5\textwidth]{Solution/AIME_2003_II_14_Solution_Figure_1.png}
                    \end{figure}
                    From this image, we can see that the y-coordinate of F is 4, and from this, we can gather that the coordinates of E, D, and C, respectively, are 8, 10, and 6.

 
                    \begin{figure}[H]
                        \centering
                        \includegraphics[width=0.5\textwidth]{Solution/AIME_2003_II_14_Solution_Figure_1.png}
                    \end{figure}
                    

 Let the angle between the $x$-axis and segment $AB$ be  $\theta$, as shown above. Thus, as $\angle FAB=120^\circ$, the angle between the $x$-axis and segment $AF$ is $60-\theta$, so $\sin{(60-\theta)}=2\sin{\theta}$. Expanding, we have

 $\sin{60}\cos{\theta}-\cos{60}\sin{\theta}=\frac{\sqrt{3}\cos{\theta}}{2}-\frac{\sin{\theta}}{2}=2\sin{\theta}$Isolating $\sin{\theta}$ we see that $\frac{\sqrt{3}\cos{\theta}}{2}=\frac{5\sin{\theta}}{2}$, or $\cos{\theta}=\frac{5}{\sqrt{3}}\sin{\theta}$. Using the fact that $\sin^2{\theta}+\cos^2{\theta}=1$, we have $\frac{28}{3}\sin^2{\theta}=1$, or $\sin{\theta}=\sqrt{\frac{3}{28}}$. Letting the side length of the hexagon be $y$, we have $\frac{2}{y}=\sqrt{\frac{3}{28}}$. After simplification we find that that $y=\frac{4\sqrt{21}}{3}$. 

 In particular, note that by the Pythagorean theorem, $b^2+2^2=y^2$, hence $b=10\sqrt{3}/3$. Also, if $C=(c,6)$, then $y^2=BC^2=4^2+(c-b)^2$, hence $c-b=8\sqrt{3}/3,$ and thus $c=18\sqrt{3}/3$. Using similar methods (or symmetry), we determine that $D=(10\sqrt{3}/3,10)$, $E=(0,8)$, and $F=(-8\sqrt{3}/3,4)$. By the Shoelace theorem,
 \[[ABCDEF]=\frac12\left|\begin{array}{cc} 0&0\\ 10\sqrt{3}/3&2\\ 18\sqrt{3}/3&6\\ 10\sqrt{3}/3&10\\ 0&8\\ -8\sqrt{3}/3&4\\ 0&0\\ \end{array}\right|=\frac12|60+180+80-36-60-(-64)|\sqrt{3}/3=48\sqrt{3}.\]

 Hence the answer is $\boxed{51}$.

 \textbf{Note} By symmetry the area of $ABCDEF$ is twice the area of $ABCF$. Therefore, you only need to calculate the coordinates of $B$, $C$, and $F$. 

 





\textbf{Solution 3}

This is similar to solution 2 but faster and easier.First off we see that the y coordinate of F must be 4, the y coordinate of E must be 8, the y coordinate of D must be 10, and the y coordinate of C must be 6 (from the parallel sides of the hexagon).We then use the sine sum angle formula*** to find the x coordinate of B (lets call it $x$):  $2\cdot\cos(120)+x\sin(120)=\frac{x\sqrt3}{2}-1=4\rightarrow x=\frac{10\sqrt3}{3}$.Now that we know $x$ we can find the x coordinate of F in multiple ways, including using the cosine sum angle formula or using the fact that triangle AFE is isosceles and AE is on the y axis.  Either way, we find that the x coordinate of F is $-\frac{8\sqrt3}{3}$.Now, divide ABCDEF into two congruent triangles and a parallelogram:  AFE, BCD, and ABDE.  The areas of AFE and BCD are each $\frac12\cdot\frac{8\sqrt3}{3}\cdot8=\frac{32\sqrt3}{3}$.  The area of ABDE is $\frac12\cdot8\cdot\frac{10\sqrt3}{3}=\frac{80\sqrt3}3$.The total area of the hexagon is $2\cdot\frac{32\sqrt3}3+\frac{80\sqrt3}3=\frac{144\sqrt3}{3}=48\sqrt3\rightarrow48+3=\boxed{051}$

 
                \begin{itemize}
                    <li>
                \begin{itemize}
                    <li>
                \begin{itemize}
                    \item 
                \end{itemize}


                </li>
                \end{itemize}


                </li>
                \end{itemize}


                The sine sum angle formula is calculated by the rotation matrix in 2D. Rotating a point counterclockwise around the origin yields new coordinates $(x' = x\cos\theta - y\sin\theta, y' = x\sin\theta + y\cos\theta)$.

 





\textbf{Solution 4 (No Trig)}


                    \begin{figure}[H]
                        \centering
                        \includegraphics[width=0.5\textwidth]{Solution/AIME_2003_II_14_Solution_Figure_1.png}
                    \end{figure}
                    

 First, we see that the y-coordinates of F, E, D, and C must be $4$, $8$, and $10$, and $6$, respectively, as in the previous solutions. We can draw a rectangle around the hexagon ABCDEF and use negative space to find the area of the hexagon. If we call the distance from the foot of the perpendiculars of B and F to A $x$ and $z$, respectively, and the distance from the bottom left vertex of the rectangle to the foot of the perpendicular from B $y$. This tells us that the area of the entire rectangle is $10(x+y+z)$, since the opposite sides are parallel and thus the length of the rectangle is $4+4+2=10$. Then, if we find the area of the extra triangles and subtract, we find that the area of hexagon ABCDEF as $6x+8z+2y$. However, noticing that $x=y$, the area of ABCDEF can also be expressed as $8(x+z)$. Now we just need to find $x+z$. Since $AB=AF$ and $\angle BAF = 120$ degrees, $BF=AB\sqrt{3}$. However, we can find AB by using the Pythagorean Theorem on either of the right triangles formed by dropping perpendiculars from B and F to the x-axis (let's call them ABX and AFY).From triangle ABX we have that $AB=\sqrt{4+x^2}$, so $BF=\sqrt{3x^2+12}$. Since AB=AF, we can also form the equation $4+x^2=16+z^2$.We can also find BF by dropping a perpendicular from B to line FY and using the Pythagorean Theorem on the right triangle formed. This gives us $BF=\sqrt{4+(x+z)^2}$. Setting our two values of BF equal and substituting $x^2$ as $12+z^2$ and simplifying, we get the equation $3z^4-16z^2-1024=0$. Now we can use the quadratic formula to get that $z^2=\frac{64}{3}$ or $-18$, so $z^2=\frac{64}{3}$. Plugging this value back into the equation $x^2=12+z^2$, we get that $x^2=\frac{100}{3}$. Now we get that $x+z$ is $6\sqrt{3}$, so the area of the hexagon is $8 \cdot 6\sqrt{3}=48\sqrt{3}$, so the answer is $48+3=\boxed{051}$

 




	\item (\textit{AIME II 2003}) Let 
 \[P(x) = 24x^{24} + \sum_{j = 1}^{23}(24 - j)(x^{24 - j} + x^{24 + j}).\] Let $z_{1},z_{2},\ldots,z_{r}$ be the distinct zeros of $P(x),$ and let $z_{k}^{2} = a_{k} + b_{k}i$ for $k = 1,2,\ldots,r,$ where $a_{k}$ and $b_{k}$ are real numbers. Let


 $\sum_{k = 1}^{r}|b_{k}| = m + n\sqrt {p},$where $m, n,$ and $p$ are integers and $p$ is not divisible by the square of any prime. Find $m + n + p.$


 
 \textbf{Answer: }015



\textbf{Solution}

This can be factored as:

 
 \[P(x) = x\left( x^{23} + x^{22} + \cdots + x^2 + x + 1 \right)^2\]

 Note that $\left( x^{23} + x^{22} + \cdots + x^2 + x + 1 \right) \cdot (x-1) = x^{24} - 1$.So the roots of $x^{23} + x^{22} + \cdots + x^2 + x + 1$ are exactly all $24$-th complex roots of $1$, except for the root $x=1$. 

 Let $\omega=\cos \frac{360^\circ}{24} + i\sin \frac{360^\circ}{24}$. Then the distinct zeros of $P$ are $0,\omega,\omega^2,\dots,\omega^{23}$.w

 We can clearly ignore the root $x=0$ as it does not contribute to the value that we need to compute. 

 The squares of the other roots are $\omega^2,~\omega^4,~\dots,~\omega^{24}=1,~\omega^{26}=\omega^2,~\dots,~\omega^{46}=\omega^{22}$.

 Hence we need to compute the following sum:

 
 \[R = \sum_{k = 1}^{23} \left|\, \sin \left( k\cdot \frac{360^\circ}{12} \right) \right|\]

 Using basic properties of the sine function, we can simplify this to

 
 \[R = 4 \cdot \sum_{k = 1}^{5} \sin \left( k\cdot \frac{360^\circ}{12} \right)\]

 The five-element sum is just $\sin 30^\circ + \sin 60^\circ + \sin 90^\circ + \sin 120^\circ + \sin 150^\circ$. We know that $\sin 30^\circ = \sin 150^\circ = \frac 12$, $\sin 60^\circ = \sin 120^\circ = \frac{\sqrt 3}2$, and $\sin 90^\circ = 1$.Hence our sum evaluates to:

 
 \[R = 4 \cdot \left( 2\cdot \frac 12 + 2\cdot \frac{\sqrt 3}2 + 1 \right) = 8 + 4\sqrt 3\]

 Therefore the answer is $8+4+3 = \boxed{015}$.

 





\textbf{Solution 2}

Note that $x^k + x^{k-1} + \dots + x + 1 = \frac{x^{k+1} - 1}{x - 1}$. Our sum can be reformed as 
 \[\frac{x(x^{47} - 1) + x^2(x^{45} - 1) + \dots + x^{24}(x - 1)}{x-1}\]

 So 
 \[\frac{x^{48} + x^{47} + x^{46} + \dots + x^{25} - x^{24} - x^{23} - \dots - x}{x-1} = 0\]

 $x(x^{47} + x^{46} + \dots - x - 1) = 0$

 $x^{47} + x^{46} + \dots - x - 1 = 0$

 $x^{47} + x^{46} + \dots + x + 1 = 2(x^{23} + x^{22} + \dots + x + 1)$

 $\frac{x^{48} - 1}{x - 1} = 2\frac{x^{24} - 1}{x - 1}$

 $x^{48} - 1 - 2x^{24} + 2 = 0$

 $(x^{24} - 1)^2 = 0$

 And we can proceed as above.

 





\textbf{Solution 3}

As in Solution 1, we find that the roots of $P(x)$ we care about are the 24th roots of unity except $1$. Therefore, the squares of these roots are the 12th roots of unity. In particular, every 12th root of unity is counted twice, except for $1$, which is only counted once. 

 The possible imaginary parts of the 12th roots of unity are $0$, $\pm\frac{1}{2}$, $\pm\frac{\sqrt{3}}{2}$, and $\pm 1$. We can disregard $0$ because it doesn't affect the sum. 

 $8$ squares of roots have an imaginary part of $\pm\frac{1}{2}$, $8$ squares of roots have an imaginary part of $\pm\frac{\sqrt{3}}{2}$, and $4$ squares of roots have an imaginary part of $\pm 1$. Therefore, the sum equals $8\left(\frac{1}{2}\right) + 8\left(\frac{\sqrt{3}}{2}\right) + 4(1) = 8 + 4\sqrt{3}$.

 The answer is $8+4+3=\boxed{015}$.

 - AMBRIGGS

 




\end{enumerate}